\section{SQL Optimization}

% ------------------------------------------------------------
\begin{frame}{Why Query Optimization Matters}
  A correct query that runs in 30 seconds is useless in a web application that
  must respond in under 200\,ms. As data volumes grow, the difference between
  an optimized and an unoptimized query can be seconds vs.\ hours.

  \begin{block}{Two Levels of Optimization}
    \begin{enumerate}
      \item \textbf{Query-level}: rewrite the SQL to be more efficient --
            reduce rows processed as early as possible
      \item \textbf{Schema-level}: add or adjust \textbf{indexes} so the
            database can locate rows without scanning the entire table
    \end{enumerate}
  \end{block}

  \begin{alertblock}{The Golden Rule}
    \textbf{Measure before you optimize.} Use \texttt{EXPLAIN} /
    \texttt{EXPLAIN ANALYZE} to find the real bottleneck.
    Optimizing the wrong query wastes time and can make things worse.
  \end{alertblock}
\end{frame}

% ------------------------------------------------------------
\begin{frame}{How a Query is Executed -- The Query Planner}
  When you send a \texttt{SELECT} statement, the DBMS does not execute it
  literally. A \textbf{query planner} produces an \emph{execution plan} --
  a tree of physical operations that produce the same result as efficiently
  as possible.

  \begin{block}{Planning Steps}
    \begin{enumerate}
      \item \textbf{Parsing}: validate syntax, resolve table/column names
      \item \textbf{Logical rewriting}: push predicates down, eliminate
            redundancy, reorder joins
      \item \textbf{Statistics lookup}: uses row count, cardinality,
            histograms to estimate costs
      \item \textbf{Physical plan selection}: choose full scan, index scan,
            nested-loop join, hash join, etc.
      \item \textbf{Execution}: run the chosen plan and stream results
    \end{enumerate}
  \end{block}

  \begin{examples}{Key Insight}
    The planner can only choose from \emph{access paths that exist}.
    Adding an index creates a new, cheaper path for the planner to choose.
  \end{examples}
\end{frame}

% ------------------------------------------------------------
\begin{frame}[fragile]{EXPLAIN -- Reading the Execution Plan}
  \begin{block}{Syntax}
\begin{lstlisting}[style=sqlstyle]
-- Show estimated plan (no actual execution)
EXPLAIN
SELECT * FROM Track WHERE AlbumID = 1;

-- Execute AND show actual timing + row counts
EXPLAIN ANALYZE
SELECT ar.Name, al.Title
FROM   Artist ar
JOIN   Album  al ON ar.ArtistID = al.ArtistID
WHERE  al.Genre = 'Electronic';
\end{lstlisting}
  \end{block}

  \begin{block}{What to Look For}
    \begin{itemize}
      \item \textbf{type = ALL}: full table scan -- red flag for large tables
      \item \textbf{type = ref / range}: index used -- good
      \item \textbf{rows}: estimated rows examined. High value = slow query
      \item \textbf{Extra: Using filesort}: sorting without index -- add one
      \item \textbf{Extra: Using temporary}: temp table created -- expensive
    \end{itemize}
  \end{block}
\end{frame}

% ------------------------------------------------------------
\begin{frame}[fragile]{EXPLAIN -- Example Output Comparison}
  \begin{block}{Query: \texttt{SELECT * FROM Track WHERE AlbumID = 3}}
\begin{lstlisting}[style=sqlstyle]
-- WITHOUT index on AlbumID:
-- type=ALL  key=NULL  rows=5  Extra=Using where
-- Reads ALL 5 rows to find matches (full table scan)

-- After: CREATE INDEX idx_track_album ON Track(AlbumID);

-- WITH index on AlbumID:
-- type=ref  key=idx_track_album  rows=2
-- Only reads 2 matching rows directly via index
\end{lstlisting}
  \end{block}

  \begin{block}{Reading the key EXPLAIN columns}
    \begin{tabular}{@{}ll@{}}
      \toprule
      \textbf{Column} & \textbf{Meaning} \\
      \midrule
      \texttt{type}   & Access method: ALL (bad) $\to$ ref/range (good) $\to$ const (best) \\
      \texttt{key}    & Index used (NULL = no index) \\
      \texttt{rows}   & Estimated rows to examine \\
      \texttt{Extra}  & Additional info: filesort, Using index (covering), etc. \\
      \bottomrule
    \end{tabular}
  \end{block}
\end{frame}
